% !TEX program = xelatex
\documentclass[12pt,a4paper]{article}
\usepackage{fontspec}
\usepackage{polyglossia}
\setmainlanguage{russian}
\setotherlanguage{english}
\setmainfont{Times New Roman}
\usepackage{amsmath,amssymb,amsfonts,mathtools}
\usepackage{geometry}
\geometry{left=2cm,right=2cm,top=2cm,bottom=2cm}
\usepackage{enumitem}
\usepackage{hyperref}
\usepackage{indentfirst}
\usepackage{titlesec}
\setlength{\parindent}{1.25cm}
\setlength{\parskip}{0.35em}
\sloppy
\title{Зачет по вычислительным методам: ответы на вопросы}
\author{}
\date{}
\begin{document}
\maketitle
\tableofcontents
\newpage

\section{Пререквизиты и основные термины}
\textbf{Стационарная точка} - это точка, в которой градиент равен нулю: $\nabla f(x^*)=0.$
В одномерном случае это условие записывается так: $f'(x^*)=0.$
Стационарная точка может быть минимумом, максимумом или седловой точкой.
\textbf{Матрица Якоби} - это матрица первых частных производных вектор-функции. Если $r(x)=(r_1(x),r_2(x),\ldots,r_m(x))^T$, где $x=(x_1,x_2,\ldots,x_n)^T$, то матрица Якоби имеет вид:
\[
J(x)=
\begin{pmatrix}
\frac{\partial r_1}{\partial x_1} &
\frac{\partial r_1}{\partial x_2} &
\cdots &
\frac{\partial r_1}{\partial x_n}\\
\frac{\partial r_2}{\partial x_1} &
\frac{\partial r_2}{\partial x_2} &
\cdots &
\frac{\partial r_2}{\partial x_n}\\
\vdots & \vdots & \ddots & \vdots\\
\frac{\partial r_m}{\partial x_1} &
\frac{\partial r_m}{\partial x_2} &
\cdots &
\frac{\partial r_m}{\partial x_n}
\end{pmatrix}.
\]
Она показывает, как каждая компонента вектор-функции меняется при изменении каждого параметра. В методе Левенберга-Марквардта матрица Якоби строится для вектора невязок $r(x)$.
\textbf{Матрица Гессе} - это матрица вторых частных производных:
\[
H(x)=\nabla^2 f(x)=
\begin{pmatrix}
\frac{\partial^2 f}{\partial x_1^2} &
\frac{\partial^2 f}{\partial x_1\partial x_2} &
\cdots &
\frac{\partial^2 f}{\partial x_1\partial x_n}\\
\frac{\partial^2 f}{\partial x_2\partial x_1} &
\frac{\partial^2 f}{\partial x_2^2} &
\cdots &
\frac{\partial^2 f}{\partial x_2\partial x_n}\\
\vdots & \vdots & \ddots & \vdots\\
\frac{\partial^2 f}{\partial x_n\partial x_1} &
\frac{\partial^2 f}{\partial x_n\partial x_2} &
\cdots &
\frac{\partial^2 f}{\partial x_n^2}
\end{pmatrix}.
\]
Она описывает кривизну функции.

\textbf{Положительно определенная матрица} - матрица $A$, для которой для любого ненулевого вектора $d$: $d^TAd>0.$
Для функции это означает: если положительно определена матрица Гессе $H(x^*)$ в стационарной точке $x^*$ , то в малой окрестности этой точки функция похожа на чашу (если около точки $x^*$ функция при любом малом сдвиге $d > 0$ увеличивается), и $x^*$ является точкой строгого локального минимума
\textbf{Отрицательно определенная матрица} - матрица $A$, для которой для любого ненулевого $d$: $d^TAd<0.$
Для функции это означает, что около точки поверхность похожа на перевернутую чашу, то есть возможен максимум.
Если выражение $d^TAd$ может быть и положительным, и отрицательным при разных $d$, матрица называется неопределенной. Тогда обычно возникает седловая точка.
\textbf{Седловая точка} - это стационарная точка, которая не является ни минимумом, ни максимумом. В одних направлениях функция возрастает, а в других убывает.
Пример: $f(x,y)=x^2-y^2.$
В точке $(0,0)$ градиент равен нулю, но минимума или максимума нет.
\textbf{Вычислительный граф} - это представление вычисления функции в виде последовательности элементарных операций. Узлы графа - промежуточные значения, ребра показывают зависимости между ними.
Например, для $f(x)=x^2+\sin x$
можно ввести промежуточные переменные: $u=x^2,\quad v=\sin x,\quad f=u+v.$
Такое представление используется в автоматическом дифференцировании.
\textbf{Правило цепочки} - позволяет находить производную сложной функции. Если $y=f(g(x)),$
то $y'=f'(g(x))g'(x).$
Для автоматического дифференцирования это главное правило распространения производных по вычислительному графу.
\textbf{Численное дифференцирование} - это приближенное вычисление производной по значениям функции в соседних точках.
Например: $f'(x)\approx \frac{f(x+h)-f(x)}{h}.$
Здесь $h$ - шаг дифференцирования.
Функция одной переменной называется \textbf{унимодальной} на отрезке $[a,b]$, если на этом отрезке она имеет один минимум: сначала убывает, а потом возрастает.
Для таких функций хорошо работают методы сужения отрезка: золотое сечение, метод Фибоначчи, деление отрезка пополам.
\textbf{Интервал неопределенности} - это отрезок, про который известно, что внутри него находится минимум. В процессе одномерной оптимизации этот интервал постепенно уменьшают.
\textbf{Константа Липшица} $L$ ограничивает скорость изменения функции: $|f(x_1)-f(x_2)|\le L|x_1-x_2|.$
Если такая константа существует, функция называется липшицевой. Это свойство используется в методе Пиявского для поиска глобального минимума.
\textbf{Точка экстремума} - точка, в которой функция достигает локального минимума или локального максимума. Точка $x^*$ называется точкой локального минимума, если существует такая окрестность $U(x^*)$, что для всех $x \in U(x^*)$ выполняется $f(x^*) \le f(x).$ Точка $x^*$ называется точкой локального максимума, если в некоторой окрестности $U(x^*)$ для всех $x$ выполняется $f(x^*) \ge f(x).$
\section{1. Необходимые и достаточные условия экстремумов функций}
Пусть дана функция $f(x)$, где $x \in \mathbb{R}^n$. Точка $x^*$ называется точкой локального минимума, если существует окрестность $U(x^*)$, такая что для всех $x \in U(x^*)$: $f(x^*) \le f(x).$
Аналогично, $x^*$ является точкой локального максимума, если $f(x^*) \ge f(x).$
\textbf{Необходимое условие первого порядка} - Если функция $f$ дифференцируема в точке $x^*$, и $x^*$ является внутренней точкой локального экстремума, то $\nabla f(x^*) = 0.$
То есть все первые частные производные равны нулю: $\frac{\partial f}{\partial x_1}(x^*) = 0,\quad \frac{\partial f}{\partial x_2}(x^*) = 0,\quad \ldots,\quad \frac{\partial f}{\partial x_n}(x^*) = 0.$
Такая точка называется стационарной. Важно: стационарная точка не обязательно является экстремумом.
\textbf{Необходимые условия второго порядка} - Если $f$ дважды дифференцируема и $x^*$ является точкой локального минимума, то матрица Гессе $H(x^*) = \nabla^2 f(x^*)$
должна быть неотрицательно определенной: $d^T H(x^*) d \ge 0$
для любого направления $d$.
Для локального максимума матрица Гессе должна быть неположительно определенной: $d^T H(x^*) d \le 0.$
\textbf{Достаточные условия второго порядка} - Если $\nabla f(x^*) = 0$
и матрица Гессе $H(x^*)$ положительно определена, то $x^*$ является точкой строгого локального минимума.
Если $\nabla f(x^*) = 0$
и $H(x^*)$ отрицательно определена, то $x^*$ является точкой строгого локального максимума.
Если матрица Гессе неопределенная, то в точке $x^*$ находится седловая точка.
\subsection{Пример}
Рассмотрим функцию $f(x,y) = x^2 + 4y^2 - 4x + 8y + 10.$
Найдем градиент:
\[
\nabla f(x,y) =
\begin{pmatrix}
2x - 4 \\
8y + 8
\end{pmatrix}.
\]
Приравняем его к нулю: $2x - 4 = 0,\quad 8y + 8 = 0.$
Отсюда $x = 2,\quad y = -1.$
Матрица Гессе:
\[
H =
\begin{pmatrix}
2 & 0 \\
0 & 8
\end{pmatrix}.
\]
Она положительно определена, так как ее собственные значения $2$ и $8$ положительны. Значит, точка $(2,-1)$
является точкой строгого локального минимума.
Значение функции: $f(2,-1) = 4 + 4 - 8 - 8 + 10 = 2.$
Итак, минимум равен $2$.
\section{2. Автоматическое дифференцирование: принципы работы и конкретные алгоритмы}
Автоматическое дифференцирование - это способ точного вычисления производных функции, заданной программой. Оно не является ни символьным дифференцированием, ни численным дифференцированием.
Основная идея: любая программа вычисления функции состоит из элементарных операций: $+,\ -, \cdot,\ /,\ \sin,\ \cos,\ \exp,\ \log$
и для каждой такой операции известна производная. Автоматическое дифференцирование применяет правило цепочки к последовательности вычислений.
\textbf{Отличие от численного дифференцирования} - При численном дифференцировании производная приближается разностной формулой: $f'(x) \approx \frac{f(x+h)-f(x)}{h}.$
Здесь возникает ошибка округления и ошибка аппроксимации.
При автоматическом дифференцировании производная вычисляется вместе со значением функции с машинной точностью, потому что используются точные правила дифференцирования элементарных операций.
\textbf{Прямой режим автоматического дифференцирования} - В прямом режиме вместе с каждой переменной хранят пару: $(v, \dot v),$
где $v$ - значение переменной, а $\dot v$ - производная этой переменной по выбранному входному аргументу.
Например: $z = x y.$
Тогда $\dot z = \dot x y + x \dot y.$
Для функции многих переменных прямой режим удобно использовать, когда входов мало, а выходов много.
\textbf{Обратный режим автоматического дифференцирования} - В обратном режиме сначала выполняется прямой проход: вычисляются и запоминаются все промежуточные значения.
Затем выполняется обратный проход: от результата к входным переменным распространяются производные по правилу цепочки.
Обратный режим особенно эффективен для функций вида $f:\mathbb{R}^n \to \mathbb{R},$
то есть когда входов много, а выход один. Поэтому он широко используется при обучении нейронных сетей.
\subsection{Конкретный алгоритм прямого режима}
Пусть надо вычислить $f(x)$ и $f'(x)$.
\noindent\textbf{1.}\ Входу $x$ сопоставляют пару $(x,1)$.\par
\noindent\textbf{2.}\ Константам сопоставляют пары $(c,0)$.\par
\noindent\textbf{3.}\ Для каждой операции пересчитывают значение и производную.\par
\noindent\textbf{4.}\ На выходе получают пару $(f(x), f'(x))$.\par
\subsection{Конкретный алгоритм обратного режима}
\noindent\textbf{1.}\ Построить вычислительный граф функции.\par
\noindent\textbf{2.}\ Выполнить прямой проход и сохранить значения всех узлов.\par
\noindent\textbf{3.}\ Производной выхода по самому себе присвоить значение $1$.\par
\noindent\textbf{4.}\ В обратном порядке обойти граф.\par
\noindent\textbf{5.}\ Для каждого узла добавить вклад в производные его аргументов.\par
\noindent\textbf{6.}\ Получить градиент по входным переменным.\par
\subsection{Пример}
Рассмотрим функцию $f(x) = x^2 + \sin x.$
Вычислим значение и производную в точке $x=0$ прямым режимом.
Вход: $x = (0,1).$
Операция $u=x^2$: $u = (0^2, 2x \cdot 1) = (0,0).$
Операция $v=\sin x$: $v = (\sin 0, \cos 0 \cdot 1) = (0,1).$
Операция $f=u+v$: $f = (0+0, 0+1) = (0,1).$
Значит, $f(0)=0,\quad f'(0)=1.$
Проверим аналитически: $f'(x)=2x+\cos x,\quad f'(0)=1.$
Результат совпадает.
\section{3. Численное дифференцирование: принципы работы и конкретные алгоритмы}
Численное дифференцирование - это приближенное вычисление производных функции по значениям функции в нескольких точках.
Оно используется, когда аналитическая производная неизвестна, сложна или функция задана как "черный ящик".
Основная идея основана на разложении функции в ряд Тейлора около точки $x$. Для значения справа от точки:
\[
f(x+h)=f(x)+h f'(x)+\frac{h^2}{2}f''(x)+\frac{h^3}{6}f'''(x)+O(h^4).
\]
Для значения слева от точки:
\[
f(x-h)=f(x)-h f'(x)+\frac{h^2}{2}f''(x)-\frac{h^3}{6}f'''(x)+O(h^4).
\]
Из этих двух разложений получают основные разностные формулы.
\textbf{Правая разностная производная} - из разложения $f(x+h)$ переносим $f(x)$ влево:
\[
f(x+h)-f(x)=h f'(x)+\frac{h^2}{2}f''(x)+O(h^3).
\]
Делим на $h$:
\[
\frac{f(x+h)-f(x)}{h}=f'(x)+\frac{h}{2}f''(x)+O(h^2).
\]
Отсюда:
\[
f'(x) \approx \frac{f(x+h)-f(x)}{h}.
\]
Главный отброшенный член пропорционален $h$, поэтому погрешность этой формулы имеет порядок $O(h)$. Формула простая, но не очень точная.
\textbf{Левая разностная производная} - аналогично используем разложение $f(x-h)$:
\[
f(x)-f(x-h)=h f'(x)-\frac{h^2}{2}f''(x)+O(h^3).
\]
Делим на $h$:
\[
\frac{f(x)-f(x-h)}{h}=f'(x)-\frac{h}{2}f''(x)+O(h^2).
\]
Отсюда:
\[
f'(x) \approx \frac{f(x)-f(x-h)}{h}.
\]
Погрешность также порядка $O(h)$.
\textbf{Центральная разностная производная} - вычитаем разложение $f(x-h)$ из разложения $f(x+h)$:
\[
f(x+h)-f(x-h)=2h f'(x)+\frac{h^3}{3}f'''(x)+O(h^5).
\]
Делим на $2h$:
\[
\frac{f(x+h)-f(x-h)}{2h}=f'(x)+\frac{h^2}{6}f'''(x)+O(h^4).
\]
Отсюда:
\[
f'(x) \approx \frac{f(x+h)-f(x-h)}{2h}.
\]
Здесь члены с $f''(x)$ взаимно уничтожаются, поэтому погрешность имеет порядок $O(h^2)$. Поэтому центральная разность обычно точнее правой и левой.
\textbf{Вторая производная} - складываем разложения $f(x+h)$ и $f(x-h)$:
\[
f(x+h)+f(x-h)=2f(x)+h^2 f''(x)+\frac{h^4}{12}f^{(4)}(x)+O(h^6).
\]
Переносим $2f(x)$ влево и делим на $h^2$:
\[
\frac{f(x+h)-2f(x)+f(x-h)}{h^2}=f''(x)+\frac{h^2}{12}f^{(4)}(x)+O(h^4).
\]
Отсюда:
\[
f''(x) \approx \frac{f(x+h)-2f(x)+f(x-h)}{h^2}.
\]
Погрешность этой формулы порядка $O(h^2)$.
\textbf{Выбор шага} - Если $h$ слишком большой, возникает большая ошибка аппроксимации. Если $h$ слишком маленький, усиливаются ошибки округления из-за вычитания близких чисел.
Поэтому на практике шаг выбирают не бесконечно малым, а умеренно малым. Для вычислений с двойной точностью часто используют величины порядка $h \sim 10^{-5}$
для первой производной.
\subsection{Конкретный алгоритм}
Для вычисления $f'(x)$ центральной разностью:
\noindent\textbf{1.}\ Выбрать шаг $h$.\par
\noindent\textbf{2.}\ Вычислить $f(x+h)$.\par
\noindent\textbf{3.}\ Вычислить $f(x-h)$.\par
\noindent\textbf{4.}\ Посчитать $f'(x) \approx \frac{f(x+h)-f(x-h)}{2h}.$\par
\subsection{Пример}
Пусть $f(x)=e^x.$
Найдем приближенно производную в точке $x=0$ центральной разностью при $h=0.1$. $f'(0) \approx \frac{f(0.1)-f(-0.1)}{2\cdot 0.1}.$
Так как $e^{0.1} \approx 1.10517,\quad e^{-0.1} \approx 0.90484,$
получаем $f'(0) \approx \frac{1.10517 - 0.90484}{0.2} = \frac{0.20033}{0.2} = 1.00165.$
Точное значение: $(e^x)' = e^x,\quad f'(0)=1.$
Ошибка составляет примерно $0.00165.$
\section{4. Численные методы оптимизации одномерных функций: метод Фибоначчи, золотого сечения, деления отрезка пополам. Способ нахождения начального интервала}
Рассматривается задача одномерной минимизации: $\min_{x \in [a,b]} f(x).$
Обычно предполагают, что функция унимодальна на отрезке $[a,b]$, то есть имеет один минимум: сначала убывает, затем возрастает.
\textbf{Нахождение начального интервала неопределенности} - Если начальный отрезок неизвестен, его можно найти поиском с расширением шага.
Алгоритм:
\noindent\textbf{1.}\ Выбрать начальную точку $x_0$ и шаг $h>0$.\par
\noindent\textbf{2.}\ Вычислить $f(x_0)$ и $f(x_0+h)$.\par
\noindent\textbf{3.}\ Если $f(x_0+h)>f(x_0)$, попробовать направление $-h$.\par
\noindent\textbf{4.}\ Двигаться в сторону убывания функции.\par
\noindent\textbf{5.}\ На каждом шаге увеличивать шаг, например в 2 раза: $h,\ 2h,\ 4h,\ 8h,\ldots$\par
\noindent\textbf{6.}\ Остановиться, когда значение функции начнет возрастать.\par
Полученные три точки $x_{k-1}, x_k, x_{k+1}$, для которых $f(x_k) < f(x_{k-1}),\quad f(x_k)<f(x_{k+1}),$
задают интервал, содержащий минимум: $[x_{k-1},x_{k+1}].$
\textbf{Метод деления отрезка пополам} - метод, в котором берут середину $c=\frac{a+b}{2}$
и две близкие точки: $x_1=c-\varepsilon,\quad x_2=c+\varepsilon.$
Если $f(x_1)<f(x_2),$
то минимум находится левее, и новый отрезок: $[a,x_2].$
Если $f(x_1)>f(x_2),$
то минимум находится правее, и новый отрезок: $[x_1,b].$
Метод прост, но требует по два вычисления функции почти на каждой итерации.
\textbf{Метод золотого сечения} - уменьшает интервал так, чтобы на каждой итерации можно было использовать одно уже вычисленное значение функции.
Используется число $\tau = \frac{\sqrt{5}-1}{2} \approx 0.618.$
Внутренние точки: $x_1=b-\tau(b-a),$ $x_2=a+\tau(b-a).$
Если $f(x_1)>f(x_2),$
то минимум лежит на $[x_1,b]$.
Если $f(x_1)\le f(x_2),$
то минимум лежит на $[a,x_2]$.
Процесс продолжают, пока длина отрезка не станет меньше заданной точности: $b-a < \varepsilon.$
\textbf{Метод Фибоначчи} - похож на метод золотого сечения, но использует числа Фибоначчи: $F_0=0,\quad F_1=1,\quad F_{k+1}=F_k+F_{k-1}.$
Если известно заранее, что нужно сделать $N$ итераций, то точки выбирают через отношения чисел Фибоначчи: $x_1 = a + \frac{F_{N-2}}{F_N}(b-a),$ $x_2 = a + \frac{F_{N-1}}{F_N}(b-a).$
После сравнения $f(x_1)$ и $f(x_2)$ отрезок сокращается. Метод Фибоначчи оптимален по числу вычислений функции при заранее заданном числе итераций.
\subsection{Пример}
Найдем минимум функции $f(x)=(x-3)^2$
на отрезке $[0,5]$ методом золотого сечения. $\tau \approx 0.618.$
Первая итерация: $x_1=5-0.618(5-0)=1.91,$ $x_2=0+0.618(5-0)=3.09.$
Значения функции: $f(1.91)=(1.91-3)^2=1.1881,$ $f(3.09)=(3.09-3)^2=0.0081.$
Так как $f(x_1)>f(x_2),$
минимум находится на отрезке $[1.91,5].$
Следующие итерации будут дальше сужать интервал вокруг точки $x=3$. Точный минимум: $x^*=3,\quad f(x^*)=0.$
\section{5. Метод квадратичных приближений, метод отыскания глобального оптимума Липшицевых функций, метод Пиявского}
\textbf{Метод квадратичных приближений} - используется для одномерной минимизации функции. Идея: на каждой итерации заменить исходную функцию квадратичной параболой, построенной по нескольким точкам, и взять вершину этой параболы как новое приближение к минимуму.
Пусть известны три точки: $x_1 < x_2 < x_3$
и значения функции: $f_1=f(x_1),\quad f_2=f(x_2),\quad f_3=f(x_3).$
Строится квадратичная интерполяция: $q(x) \approx f(x).$
Минимум параболы берется как новая точка: $\bar{x} = x_2 - \frac{1}{2} \frac{(x_2-x_1)^2(f_2-f_3)-(x_2-x_3)^2(f_2-f_1)} {(x_2-x_1)(f_2-f_3)-(x_2-x_3)(f_2-f_1)}.$
Затем из четырех точек выбирают три, которые лучше окружают минимум, и процесс повторяют.
Метод хорошо работает, если функция гладкая и около минимума похожа на параболу.
\subsection{Пример метода квадратичных приближений}
Пусть $f(x)=(x-2)^2+1.$
Возьмем точки $x_1=0,\quad x_2=1,\quad x_3=4.$
Значения функции: $f(0)=5,\quad f(1)=2,\quad f(4)=5.$
Поскольку функция сама является параболой, квадратичное приближение совпадает с ней точно. Вершина параболы находится в точке $x=2.$
Следовательно, метод сразу находит точный минимум: $x^*=2,\quad f(x^*)=1.$
\textbf{Липшицева функция} - функция $f(x)$ на отрезке $[a,b]$, для которой существует число $L>0$, такое что для любых $x_1,x_2 \in [a,b]$: $|f(x_1)-f(x_2)| \le L |x_1-x_2|.$
Число $L$ называется константой Липшица.
Из этого условия следует, что функция не может изменяться слишком резко. Это позволяет строить нижние оценки функции и искать глобальный минимум.
\textbf{Метод Пиявского} - предназначен для поиска глобального минимума липшицевой функции на отрезке $[a,b]$.
Задача: $\min_{x \in [a,b]} f(x).$
Пусть известна константа Липшица $L$. Для каждой уже вычисленной точки $x_i$ можно построить нижнюю оценку: $f(x) \ge f(x_i) - L|x-x_i|.$
График этой оценки имеет вид перевернутого конуса. Нижняя огибающая всех таких конусов дает глобальную нижнюю оценку функции.
Алгоритм метода Пиявского:
\noindent\textbf{1.}\ Задать отрезок $[a,b]$, точность $\varepsilon$ и константу Липшица $L$.\par
\noindent\textbf{2.}\ Вычислить $f(a)$ и $f(b)$.\par
\noindent\textbf{3.}\ Для каждого интервала $[x_i,x_{i+1}]$ вычислить точку, где пересекаются нижние оценки от концов: $x = \frac{x_i+x_{i+1}}{2} - \frac{f(x_{i+1})-f(x_i)}{2L}.$\par
\noindent\textbf{4.}\ Для этой точки вычислить нижнюю оценку: $R_i = \frac{f(x_i)+f(x_{i+1})}{2} - \frac{L}{2}(x_{i+1}-x_i).$\par
\noindent\textbf{5.}\ Выбрать интервал с наименьшим $R_i$, потому что там потенциально может находиться глобальный минимум.\par
\noindent\textbf{6.}\ Вычислить функцию в новой точке.\par
\noindent\textbf{7.}\ Добавить точку в список и повторять до достижения точности.\par
Метод Пиявского относится к глобальным методам: он не просто уточняет ближайший локальный минимум, а стремится доказательно не пропустить лучший минимум на всем отрезке.
\subsection{Пример метода Пиявского}
Рассмотрим функцию $f(x)=|x-1|+0.5x$
на отрезке $[0,3]$.
Эта функция липшицева. Можно взять, например, $L=1.5,$
потому что наклон функции по модулю не превосходит $1.5$.
Вычислим значения на концах: $f(0)=1,$ $f(3)=|3-1|+1.5=3.5.$
Новая точка по формуле Пиявского: $x = \frac{0+3}{2} - \frac{3.5-1}{2\cdot 1.5} = 1.5 - \frac{2.5}{3} \approx 0.667.$
Вычислим значение функции: $f(0.667)=|0.667-1|+0.5\cdot 0.667 \approx 0.333+0.333 =0.666.$
Это меньше, чем значения на концах, поэтому дальнейший поиск будет уточнять минимум вблизи точки $x=1$. Точно:
если $x<1$, то $f(x)=1-x+0.5x=1-0.5x,$
и функция убывает.
если $x>1$, то $f(x)=x-1+0.5x=1.5x-1,$
и функция возрастает.
Значит, глобальный минимум находится в точке $x^*=1,\quad f(1)=0.5.$
\section{6. Градиентный спуск и его вариации: базовый, с инерцией, с инерцией Нестерова}
Градиентный спуск - это итерационный метод минимизации функции. Рассматривается задача $\min_x f(x).$
Идея метода: из текущей точки $x_k$ перейти в сторону наибольшего убывания функции. Так как градиент $\nabla f(x_k)$ показывает направление наибольшего возрастания, для минимизации используют противоположное направление: $-\nabla f(x_k).$
\textbf{Базовый градиентный спуск} - базовый метод с формулой: $x_{k+1}=x_k-\alpha_k \nabla f(x_k),$
где:
\noindent\textbullet\ $x_k$ - текущее приближение;\par
\noindent\textbullet\ $\alpha_k>0$ - шаг обучения или длина шага;\par
\noindent\textbullet\ $\nabla f(x_k)$ - градиент функции в точке $x_k$.\par
Если шаг слишком большой, метод может перескакивать через минимум и расходиться. Если шаг слишком маленький, сходимость будет медленной.
\textbf{Пример базового градиентного спуска}
Пусть $f(x)=x^2.$
Тогда $f'(x)=2x.$
Возьмем $x_0=4$ и $\alpha=0.1$.
Первая итерация: $x_1=x_0-\alpha f'(x_0)=4-0.1\cdot 8=3.2.$
Вторая итерация: $x_2=3.2-0.1\cdot 6.4=2.56.$
Последовательность приближается к точке минимума: $x^*=0.$
\textbf{Градиентный спуск с инерцией} - метод, который учитывает не только текущий градиент, но и предыдущее направление движения. Это помогает быстрее двигаться по длинным пологим направлениям и уменьшать колебания.
Вводится скорость $v_k$: $v_{k+1}=\beta v_k-\alpha \nabla f(x_k),$ $x_{k+1}=x_k+v_{k+1}.$
Здесь:
\noindent\textbullet\ $v_k$ - накопленное направление движения;\par
\noindent\textbullet\ $\beta \in [0,1)$ - коэффициент инерции;\par
\noindent\textbullet\ $\alpha$ - шаг.\par
Если $\beta=0$, метод превращается в обычный градиентный спуск.
\textbf{Пример метода с инерцией}
Пусть функция имеет вытянутые линии уровня: $f(x,y)=x^2+100y^2.$
Обычный градиентный спуск может сильно колебаться по направлению $y$, потому что производная по $y$ большая:
\[
\nabla f(x,y)=
\begin{pmatrix}
2x\\
200y
\end{pmatrix}.
\]
Инерция сглаживает такие колебания: если несколько шагов подряд направление примерно одинаковое, скорость накапливается; если направление меняется туда-сюда, инерция частично компенсирует резкие изменения.
\textbf{Метод Нестерова} - отличается тем, что градиент считается не в текущей точке, а в точке с предварительным инерционным сдвигом.
Одна из форм записи: $v_{k+1}=\beta v_k-\alpha \nabla f(x_k+\beta v_k),$ $x_{k+1}=x_k+v_{k+1}.$
Смысл: метод сначала "смотрит вперед" по направлению накопленной скорости, а затем корректирует движение с учетом градиента в этой прогнозной точке.
\textbf{Пример идеи Нестерова}
Пусть текущая точка $x_k$ уже движется к минимуму с некоторой скоростью $v_k$. Обычная инерция сначала сделает шаг, а потом узнает, что могла перелететь минимум.
Метод Нестерова заранее проверяет точку $x_k+\beta v_k$
и считает градиент там. Поэтому он может раньше уменьшить скорость перед минимумом.
\section{7. Градиентный спуск и его вариации: AdaGrad, RMSProp, Adam}
Методы AdaGrad, RMSProp и Adam относятся к адаптивным вариантам градиентного спуска. Их идея: выбирать разные эффективные шаги для разных координат.
Это полезно, когда по одним координатам функция меняется резко, а по другим медленно.
\textbf{AdaGrad} - накапливает квадраты прошлых градиентов: $G_k=G_{k-1}+g_k^2,$
где $g_k=\nabla f(x_k).$
Операция $g_k^2$ понимается покоординатно.
Обновление: $x_{k+1}=x_k-\frac{\alpha}{\sqrt{G_k}+\varepsilon}g_k.$
Здесь $\varepsilon$ - маленькое число для защиты от деления на ноль.
Если по координате часто были большие градиенты, $G_k$ растет, и шаг по этой координате уменьшается.
Недостаток AdaGrad: сумма $G_k$ постоянно растет, поэтому шаги со временем могут стать слишком маленькими.
\subsection{Пример AdaGrad}
Пусть в некоторой координате градиенты часто равны примерно $10$, а в другой - примерно $0.1$.
Для первой координаты сумма квадратов быстро растет: $10^2+10^2+\ldots$
Поэтому шаг по ней уменьшается.
Для второй координаты сумма растет медленно: $0.1^2+0.1^2+\ldots$
Поэтому шаг остается относительно большим.

\textbf{RMSProp} - исправляет проблему AdaGrad: вместо полной суммы квадратов градиентов он использует экспоненциальное скользящее среднее. $s_k=\rho s_{k-1}+(1-\rho)g_k^2,$ $x_{k+1}=x_k-\frac{\alpha}{\sqrt{s_k}+\varepsilon}g_k.$
Здесь $\rho$ обычно близко к $1$, например $0.9$.
Старые градиенты постепенно забываются, поэтому шаг не обязательно стремится к нулю.
\subsection{Пример RMSProp}
Если функция сначала имела большие градиенты, а потом стала более пологой, AdaGrad все равно помнит все старые большие градиенты. RMSProp постепенно уменьшает их влияние, поэтому может продолжать делать разумные шаги.

\textbf{Adam} - объединяет две идеи:
\noindent\textbullet\ инерцию, то есть среднее значение градиента;\par
\noindent\textbullet\ адаптивный шаг, то есть среднее значение квадрата градиента.\par
Формулы: $m_k=\beta_1m_{k-1}+(1-\beta_1)g_k,$ $s_k=\beta_2s_{k-1}+(1-\beta_2)g_k^2.$
Здесь $m_k$ - оценка среднего градиента, $s_k$ - оценка среднего квадрата градиента.
Так как в начале $m_0=0$ и $s_0=0$, используют коррекцию смещения: $\hat m_k=\frac{m_k}{1-\beta_1^k},$ $\hat s_k=\frac{s_k}{1-\beta_2^k}.$
Обновление: $x_{k+1}=x_k-\alpha\frac{\hat m_k}{\sqrt{\hat s_k}+\varepsilon}.$
Обычно используют параметры: $\beta_1=0.9,\quad \beta_2=0.999.$
\subsection{Пример Adam}
В обучении нейронной сети функция ошибки зависит от большого числа параметров. Для разных параметров масштабы градиентов могут сильно отличаться.
Adam автоматически уменьшает шаги для параметров с большими и нестабильными градиентами и сохраняет более заметные шаги для параметров, где градиенты малы. Поэтому Adam часто сходится быстрее обычного градиентного спуска.
\section{8. Метод сопряженных градиентов для решения общей задачи оптимизации}
Метод сопряженных градиентов сначала был разработан для решения систем линейных уравнений $Ax=b,$
где $A$ - симметричная положительно определенная матрица.
Такая система эквивалентна задаче минимизации квадратичной функции: $f(x)=\frac12 x^TAx-b^Tx.$
Градиент этой функции: $\nabla f(x)=Ax-b.$
Решение системы $Ax=b$ соответствует условию $\nabla f(x)=0.$
\textbf{Идея сопряженных направлений} - В обычном градиентном спуске направления могут мешать друг другу: новый шаг частично портит результат предыдущего шага.
Метод сопряженных градиентов строит направления $p_k$, которые являются сопряженными относительно матрицы $A$: $p_i^T A p_j=0,\quad i\ne j.$
Это сильнее обычной ортогональности и учитывает форму квадратичной функции.
\subsection{Алгоритм для квадратичной функции}
Пусть $f(x)=\frac12 x^TAx-b^Tx.$
\noindent\textbf{1.}\ Выбрать начальное приближение $x_0$.\par
\noindent\textbf{2.}\ Посчитать остаток: $r_0=b-Ax_0.$\par
\noindent\textbf{3.}\ Задать первое направление: $p_0=r_0.$\par
\noindent\textbf{4.}\ На каждой итерации вычислять шаг: $\alpha_k=\frac{r_k^Tr_k}{p_k^TAp_k}.$\par
\noindent\textbf{5.}\ Обновлять точку: $x_{k+1}=x_k+\alpha_k p_k.$\par
\noindent\textbf{6.}\ Обновлять остаток: $r_{k+1}=r_k-\alpha_k A p_k.$\par
\noindent\textbf{7.}\ Вычислять коэффициент: $\beta_k=\frac{r_{k+1}^Tr_{k+1}}{r_k^Tr_k}.$\par
\noindent\textbf{8.}\ Обновлять направление: $p_{k+1}=r_{k+1}+\beta_k p_k.$\par
Для квадратичной функции в точной арифметике метод находит минимум не более чем за $n$ итераций, где $n$ - размерность задачи.
\textbf{Обобщение на неквадратичные функции} - Для общей задачи $\min_x f(x)$
используют нелинейный метод сопряженных градиентов.
Вместо остатка берут антиградиент: $r_k=-\nabla f(x_k).$
Направление: $p_k=-\nabla f(x_k)+\beta_k p_{k-1}.$
Шаг $\alpha_k$ выбирают одномерной минимизацией: $\alpha_k=\arg\min_{\alpha>0} f(x_k+\alpha p_k).$
Один из вариантов коэффициента Флетчера-Ривса: $\beta_k^{FR}= \frac{\nabla f(x_k)^T\nabla f(x_k)} {\nabla f(x_{k-1})^T\nabla f(x_{k-1})}.$
Другой вариант - Полак-Рибьер: $\beta_k^{PR}= \frac{\nabla f(x_k)^T(\nabla f(x_k)-\nabla f(x_{k-1}))} {\nabla f(x_{k-1})^T\nabla f(x_{k-1})}.$
\subsection{Пример}
Рассмотрим квадратичную функцию $f(x,y)=x^2+2y^2.$
Ее минимум находится в точке $(0,0).$
Градиент:
\[
\nabla f(x,y)=
\begin{pmatrix}
2x\\
4y
\end{pmatrix}.
\]
Эту функцию можно записать в квадратичном виде $f(z)=\frac12 z^TAz,$ где $z=(x,y)^T$, а
\[
A=
\begin{pmatrix}
2 & 0\\
0 & 4
\end{pmatrix}.
\]
Тогда $\nabla f(z)=Az.$ Возьмем начальную точку $z_0=(2,1)^T.$
Сначала считаем градиент:
$g_0=\nabla f(z_0)=Az_0=(4,4)^T.$
Первое направление берется как антиградиент:
$p_0=-g_0=(-4,-4)^T.$
Шаг для квадратичной функции считается по формуле:
$\alpha_0=-\frac{g_0^Tp_0}{p_0^TAp_0}.$
Здесь $g_0^Tp_0=(4,4)(-4,-4)^T=-32$, а $p_0^TAp_0=(-4,-4)(-8,-16)^T=96.$
Значит:
$\alpha_0=\frac{32}{96}=\frac13.$
Обновляем точку:
$z_1=z_0+\alpha_0p_0=(2,1)^T+\frac13(-4,-4)^T=\left(\frac23,-\frac13\right)^T.$
После первого шага мы минимизировали функцию вдоль направления $p_0$, но еще не пришли в общий минимум.
Считаем новый градиент:
$g_1=Az_1=\left(\frac43,-\frac43\right)^T.$
Если идти просто по антиградиенту $-g_1$, то обычный градиентный спуск может начать зигзаг. Метод сопряженных градиентов строит новое направление как смесь антиградиента и предыдущего направления:
$p_1=-g_1+\beta_0p_0.$
Для варианта Флетчера-Ривса:
$\beta_0=\frac{g_1^Tg_1}{g_0^Tg_0}=\frac{\frac{16}{9}+\frac{16}{9}}{16+16}=\frac19.$
Тогда:
$p_1=-\left(\frac43,-\frac43\right)^T+\frac19(-4,-4)^T=\left(-\frac{16}{9},\frac89\right)^T.$
Проверим главное свойство: направления должны быть сопряженными относительно $A$, то есть $p_0^TAp_1=0.$
Сначала $Ap_1=\left(-\frac{32}{9},\frac{32}{9}\right)^T.$
Тогда:
$p_0^TAp_1=(-4,-4)\left(-\frac{32}{9},\frac{32}{9}\right)^T=\frac{128}{9}-\frac{128}{9}=0.$
Значит, новое направление действительно не разрушает минимум, найденный вдоль предыдущего направления.
Считаем второй шаг:
$\alpha_1=-\frac{g_1^Tp_1}{p_1^TAp_1}.$
Здесь $g_1^Tp_1=-\frac{32}{9}$, а $p_1^TAp_1=\frac{256}{27}$.
Поэтому:
$\alpha_1=\frac{32}{9}\cdot\frac{27}{256}=\frac38.$
Обновляем точку:
$z_2=z_1+\alpha_1p_1=\left(\frac23,-\frac13\right)^T+\frac38\left(-\frac{16}{9},\frac89\right)^T=(0,0)^T.$
За две итерации метод пришел в точный минимум, потому что задача двумерная и функция квадратичная. Обычный градиентный спуск на такой вытянутой функции мог бы идти зигзагом, а метод сопряженных градиентов выбирает направления так, чтобы они были $A$-сопряженными и не портили результат предыдущих шагов.
\section{9. Способы задания шага в градиентных методах. Условия Вольфе}
В градиентных методах общий вид итерации: $x_{k+1}=x_k+\alpha_k p_k,$
где:
\noindent\textbullet\ $p_k$ - направление спуска;\par
\noindent\textbullet\ $\alpha_k>0$ - шаг.\par
Обычно для градиентного спуска $p_k=-\nabla f(x_k).$
Качество метода сильно зависит от выбора $\alpha_k$.
\textbf{Постоянный шаг} - Самый простой вариант: $\alpha_k=\alpha.$
Преимущество - простота. Недостаток - сложно подобрать универсальное значение.
Если $\alpha$ слишком велик, метод может расходиться. Если $\alpha$ слишком мал, метод сходится медленно.
\textbf{Убывающий шаг} - Можно брать шаг, который уменьшается с номером итерации: $\alpha_k=\frac{\alpha_0}{k+1}.$
Такой подход может быть полезен в стохастических методах, но в обычной оптимизации часто сходится медленно.
\textbf{Точный одномерный поиск} - На каждой итерации можно выбирать шаг как решение одномерной задачи: $\alpha_k=\arg\min_{\alpha>0} f(x_k+\alpha p_k).$
Это дает хороший шаг, но может быть дорого, потому что на каждой итерации нужно решать отдельную задачу минимизации.
\textbf{Неточный одномерный поиск} - На практике часто достаточно выбрать шаг, который не обязательно минимален, но дает достаточное уменьшение функции. Для этого используют backtracking line search и условия Вольфе.
\textbf{Правило Армихо} - Условие достаточного убывания: $f(x_k+\alpha p_k)\le f(x_k)+c_1\alpha \nabla f(x_k)^T p_k,$
где $0<c_1<1.$
Так как $p_k$ - направление спуска, $\nabla f(x_k)^T p_k<0.$
Правая часть меньше $f(x_k)$, поэтому условие требует реального уменьшения функции.
\textbf{Условия Вольфе} - состоят из двух требований.
Первое - условие Армихо: $f(x_k+\alpha p_k)\le f(x_k)+c_1\alpha \nabla f(x_k)^T p_k.$
Второе - условие кривизны: $\nabla f(x_k+\alpha p_k)^T p_k \ge c_2\nabla f(x_k)^T p_k,$
где $0<c_1<c_2<1.$
Первое условие не дает
взять неэффективный шаг, потому что функция должна достаточно
уменьшиться
Второе условие не дает взять слишком маленький шаг, потому что требует, чтобы производная вдоль направления заметно изменилась.
\textbf{Сильные условия Вольфе} - В сильных условиях Вольфе второе условие заменяется на: $|\nabla f(x_k+\alpha p_k)^T p_k| \le c_2|\nabla f(x_k)^T p_k|.$
Они часто используются в квазиньютоновских методах и методе сопряженных градиентов.
\subsection{Пример выбора шага}
Пусть $f(x)=x^2,\quad x_0=2.$
Градиент: $f'(x)=2x,\quad f'(2)=4.$
Направление спуска: $p_0=-4.$
Итерация: $x_1=2+\alpha(-4)=2-4\alpha.$
Если взять $\alpha=0.5$, то $x_1=0,$
и минимум найден сразу.
Если взять $\alpha=1$, то $x_1=-2,$
значение функции не уменьшится: $f(-2)=4=f(2).$
Если взять $\alpha=0.1$, то $x_1=1.6,\quad f(1.6)=2.56,$
уменьшение есть, но шаг не оптимальный.
\section{10. Метод Ньютона, метод Левенберга-Марквардта}
\textbf{Метод Ньютона для оптимизации} - метод, который использует не только градиент, но и матрицу вторых производных - Гессе.
Для минимизации функции $\min_x f(x)$
в окрестности текущей точки $x_k$ строят квадратичное приближение: $f(x_k+p)\approx f(x_k)+\nabla f(x_k)^Tp+\frac12p^T H(x_k)p.$
Минимум этого квадратичного приближения находится из условия: $\nabla f(x_k)+H(x_k)p=0.$
Отсюда направление Ньютона: $p_k=-H(x_k)^{-1}\nabla f(x_k).$
Итерация: $x_{k+1}=x_k+p_k.$
Часто используют вариант с шагом: $x_{k+1}=x_k+\alpha_k p_k.$
\subsection{Достоинства и недостатки метода Ньютона}
Достоинство: около минимума метод может сходиться очень быстро, квадратично.
Недостатки:
\noindent\textbullet\ нужно вычислять матрицу Гессе;\par
\noindent\textbullet\ нужно решать систему линейных уравнений;\par
\noindent\textbullet\ если Гессе не положительно определена, направление может не быть направлением спуска;\par
\noindent\textbullet\ вдали от минимума метод может вести себя неустойчиво.\par
\subsection{Пример метода Ньютона}
Пусть $f(x)=x^4+x^2.$
Тогда $f'(x)=4x^3+2x,$ $f''(x)=12x^2+2.$
Формула Ньютона в одномерном случае: $x_{k+1}=x_k-\frac{f'(x_k)}{f''(x_k)}.$
Возьмем $x_0=1.$
Первый шаг:
\[
f'(1)=4\cdot 1^3+2\cdot 1=6,
\]
\[
f''(1)=12\cdot 1^2+2=14.
\]
Тогда:
\[
x_1=x_0-\frac{f'(x_0)}{f''(x_0)}
=1-\frac{6}{14}
=1-\frac{3}{7}
=\frac{4}{7}\approx 0.571.
\]
Второй шаг считаем уже из точки $x_1=\frac{4}{7}$:
\[
f'\left(\frac{4}{7}\right)
=4\left(\frac{4}{7}\right)^3+2\left(\frac{4}{7}\right)
=\frac{256}{343}+\frac{8}{7}
=\frac{648}{343}.
\]
\[
f''\left(\frac{4}{7}\right)
=12\left(\frac{4}{7}\right)^2+2
=\frac{192}{49}+2
=\frac{290}{49}.
\]
Тогда:
\[
x_2=x_1-\frac{f'(x_1)}{f''(x_1)}
=\frac{4}{7}-\frac{\frac{648}{343}}{\frac{290}{49}}
=\frac{4}{7}-\frac{324}{1015}
=\frac{256}{1015}
\approx 0.252.
\]
Минимум этой функции находится в точке $x=0$, потому что $f(x)=x^4+x^2\ge 0$ и $f(0)=0.$ После двух шагов метод Ньютона приблизился к минимуму: $x_0=1,\quad x_1\approx 0.571,\quad x_2\approx 0.252.$

\textbf{Метод Левенберга-Марквардта} - используется в задачах нелинейного метода наименьших квадратов: $\min_x F(x)=\frac12\sum_{i=1}^m r_i(x)^2.$
Здесь $r_i(x)$ - невязки, то есть ошибки модели. (разница между тем, что предсказала модель, и фактическим значением)
В векторной форме: $F(x)=\frac12\|r(x)\|^2.$
Метод Левенберга-Марквардта можно понимать как устойчивый вариант метода Ньютона для такой суммы квадратов. Вместо полной матрицы Гессе он использует приближение через матрицу Якоби невязок: $H(x)\approx J(x)^TJ(x).$
Чтобы шаг не был слишком резким и система лучше решалась, к этому приближению добавляют демпфирование: $(J^TJ+\lambda I)p=-J^Tr.$
Здесь:
\noindent\textbullet\ $\lambda>0$ - параметр демпфирования;\par
\noindent\textbullet\ $I$ - единичная матрица.\par
Если $\lambda$ маленькое, метод ведет себя почти как быстрый ньютоновский метод для задачи наименьших квадратов.
Если $\lambda$ большое, шаг становится более осторожным и ближе по смыслу к градиентному спуску.
\textbf{Идея выбора $\lambda$} - Если после шага значение функции уменьшилось, шаг был удачным, и $\lambda$ можно уменьшить.
Если значение функции увеличилось, шаг был неудачным, и $\lambda$ увеличивают, чтобы сделать следующий шаг осторожнее.
\textbf{Алгоритм метода Левенберга-Марквардта}:
\noindent\textbf{1.}\ Выбрать начальное приближение $x_0$, начальный параметр демпфирования $\lambda>0$ и точность $\varepsilon$.\par
\noindent\textbf{2.}\ На текущей итерации $k$ вычислить вектор невязок:\par
\[
r(x_k)=
\begin{pmatrix}
r_1(x_k)\\
r_2(x_k)\\
\vdots\\
r_m(x_k)
\end{pmatrix}.
\]
\noindent\textbf{3.}\ Вычислить матрицу Якоби невязок:\par
\[
J(x_k)=
\begin{pmatrix}
\frac{\partial r_1}{\partial x_1} & \cdots & \frac{\partial r_1}{\partial x_n}\\
\frac{\partial r_2}{\partial x_1} & \cdots & \frac{\partial r_2}{\partial x_n}\\
\vdots & \ddots & \vdots\\
\frac{\partial r_m}{\partial x_1} & \cdots & \frac{\partial r_m}{\partial x_n}
\end{pmatrix}.
\]
\noindent\textbf{4.}\ Составить и решить систему для шага $p_k$:\par
\[
\left(J(x_k)^TJ(x_k)+\lambda I\right)p_k=-J(x_k)^Tr(x_k).
\]
\noindent\textbf{5.}\ Проверить новую точку:\par
\[
x_{\text{new}}=x_k+p_k.
\]
Если $F(x_{\text{new}})<F(x_k)$, то шаг принимается:
\[
x_{k+1}=x_{\text{new}},
\]
а параметр $\lambda$ уменьшается, например:
\[
\lambda:=\frac{\lambda}{10}.
\]
\noindent\textbf{6.}\ Если $F(x_{\text{new}})\ge F(x_k)$, то шаг отклоняется, точка не меняется, а $\lambda$ увеличивается:\par
\[
\lambda:=10\lambda.
\]
После увеличения $\lambda$ систему решают заново, чтобы получить более осторожный шаг.
\noindent\textbf{7.}\ Остановиться, если шаг или градиент стали достаточно малы, например:\par
\[
\|p_k\|<\varepsilon
\]
или
\[
\|J(x_k)^Tr(x_k)\|<\varepsilon.
\]
Смысл алгоритма: при маленьком $\lambda$ метод похож на Гаусса-Ньютона и делает быстрые шаги, а при большом $\lambda$ он становится ближе к градиентному спуску и двигается осторожнее.
\subsection{Пример метода Левенберга-Марквардта}
Пусть есть точки $(1,2),(2,4),(3,6)$, и нужно подобрать параметр $a$ модели $y=ax$.
Невязки: $r_i(a)=ax_i-y_i$, целевая функция: $F(a)=\frac12\sum_i(ax_i-y_i)^2.$
Правильное значение здесь $a=2$, но начнем с $a_0=0$ и возьмем $\lambda=1$.

При $a_0=0$ получаем $r(a_0)=(-2,-4,-6)^T$, а матрица Якоби по параметру $a$ равна $J=(1,2,3)^T$, потому что $\frac{\partial r_i}{\partial a}=x_i.$
Тогда $J^TJ=1^2+2^2+3^2=14$, $J^Tr=1\cdot(-2)+2\cdot(-4)+3\cdot(-6)=-28.$
Шаг находится из уравнения $(J^TJ+\lambda I)p=-J^Tr$, то есть $(14+1)p=28.$
Отсюда $p=\frac{28}{15}\approx 1.867$, поэтому $a_1=a_0+p\approx 1.867.$

Ошибка уменьшилась: $F(a_0)=28$, а $F(a_1)\approx 0.124.$
Значит, шаг удачный, и $\lambda$ можно уменьшить, например до $0.1$.
Для второго шага $J^Tr=14a_1-28=-\frac{28}{15}$, поэтому $(14+0.1)p=\frac{28}{15}$ и $p\approx 0.132.$
Тогда $a_2=a_1+p\approx 1.999$, то есть метод почти пришел к правильному значению $a=2$.

Видно, что $\lambda$ делает шаг осторожнее: без регуляризации первый шаг был бы $p=2$, а при $\lambda=1$ получился $p\approx 1.867$.


\section{11. Алгоритмы BFGS и DFP}
Алгоритмы \textbf{BFGS} и \textbf{DFP} относятся к квазиньютоновским методам оптимизации. Они применяются для решения задачи $\min_x f(x)$ где $f(x)$ - гладкая функция многих переменных.
Идея квазиньютоновских методов состоит в том, чтобы использовать направление, похожее на ньютоновское, но не вычислять матрицу Гессе напрямую. В методе Ньютона направление задается формулой:
\[
p_k=-H(x_k)^{-1}\nabla f(x_k).
\]
В квазиньютоновских методах вместо настоящей обратной матрицы Гессе используют ее приближение $B_k$:
\[
p_k=-B_k\nabla f(x_k).
\]
После каждого шага матрица $B_k$ уточняется по информации о том, как изменился градиент.
Обозначим:
\[
s_k=x_{k+1}-x_k,\qquad y_k=\nabla f(x_{k+1})-\nabla f(x_k).
\]
Здесь $s_k$ - изменение точки, а $y_k$ - изменение градиента.
Главное условие для обновления матрицы называется \textbf{секущим условием}:
\[
B_{k+1}y_k=s_k.
\]
Оно означает, что новая матрица должна правильно отражать изменение градиента на последнем шаге.
\textbf{Метод DFP} - один из первых квазиньютоновских методов. Он обновляет приближение обратной матрицы Гессе по формуле:
\[
B_{k+1}=B_k+\frac{s_ks_k^T}{s_k^Ty_k}-\frac{B_ky_ky_k^TB_k}{y_k^TB_ky_k}.
\]
Если матрица $B_k$ была положительно определенной и выполняется условие $s_k^Ty_k>0$, то новая матрица тоже остается положительно определенной.
\textbf{Метод BFGS} - более устойчивый и чаще используемый вариант. Формула обновления:
\[
B_{k+1}=
\left(I-\rho_ks_ky_k^T\right)B_k
\left(I-\rho_ky_ks_k^T\right)
+\rho_ks_ks_k^T,
\qquad
\rho_k=\frac{1}{y_k^Ts_k}.
\]
Важно: $\rho_k$ - это скаляр, а $s_ky_k^T$ - матрица. Поэтому произведение $\rho_ks_ky_k^T$ не равно $1$. Единицей становится только скалярное произведение $\rho_k y_k^Ts_k=1$. Порядок умножения здесь важен: $y_k^Ts_k$ - число, а $s_ky_k^T$ - матрица.
Метод BFGS обычно работает лучше DFP на практике, потому что устойчивее на плохо обусловленных задачах.
Общий алгоритм:
\noindent\textbf{1.}\ Выбрать начальную точку $x_0$ и начальную матрицу $B_0$, часто $B_0=I$.\par
\noindent\textbf{2.}\ Найти направление $p_k=-B_k\nabla f(x_k)$.\par
\noindent\textbf{3.}\ Выбрать шаг $\alpha_k$ с помощью одномерного поиска.\par
\noindent\textbf{4.}\ Обновить точку:\par
\[
x_{k+1}=x_k+\alpha_kp_k.
\]
\noindent\textbf{5.}\ Посчитать $s_k$, $y_k$ и обновить матрицу $B_k$.\par
\noindent\textbf{6.}\ Повторять, пока $\|\nabla f(x_k)\|$ не станет достаточно малой.\par
\subsection{Достоинства и недостатки}
Достоинства:
\noindent\textbullet\ не нужно вычислять матрицу Гессе;\par
\noindent\textbullet\ сходимость обычно быстрее, чем у обычного градиентного спуска;\par
\noindent\textbullet\ BFGS хорошо работает на широком классе гладких задач.\par
Недостатки:
\noindent\textbullet\ нужно хранить матрицу размера $n\times n$;\par
\noindent\textbullet\ для очень больших задач используют специальные варианты, например L-BFGS;\par
\noindent\textbullet\ качество работы зависит от выбора шага $\alpha_k$.\par

\subsection{Пример}
Пусть $f(x,y)=x^2+4y^2$, тогда $\nabla f(x,y)=(2x,8y)^T.$
Возьмем начальную точку $x_0=(1,1)^T$ и $B_0=I.$
Тогда $g_0=\nabla f(x_0)=(2,8)^T$, а первое направление:
\[
p_0=-B_0g_0=(-2,-8)^T.
\]
При точном одномерном поиске для этой квадратичной функции:
\[
\alpha_0=-\frac{g_0^Tp_0}{p_0^THp_0}=\frac{68}{520}\approx 0.131,
\]
где $H=\begin{pmatrix}2&0\\0&8\end{pmatrix}.$
Получаем новую точку:
\[
x_1=x_0+\alpha_0p_0\approx (0.738,-0.046)^T.
\]
Теперь считаем изменения:
\[
s_0=x_1-x_0\approx(-0.262,-1.046)^T,\qquad
y_0=g_1-g_0\approx(-0.523,-8.369)^T.
\]
Именно по векторам $s_0$ и $y_0$ метод BFGS или DFP обновляет матрицу $B_0$, чтобы получить более точное приближение к обратной матрице Гессе.
После этого следующее направление уже считается не просто как антиградиент, а как:
\[
p_1=-B_1g_1.
\]
Поэтому метод начинает учитывать, что по оси $y$ функция имеет большую кривизну, чем по оси $x$, и выбирает направление умнее обычного градиентного спуска.

\section{12. Координатная оптимизация, метод Пауэлла, Хуке-Дживса}
\textbf{Координатная оптимизация} - это метод, в котором многомерная задача оптимизации сводится к последовательности одномерных задач. На каждом шаге меняют только одну координату, а остальные фиксируют.
Пусть нужно минимизировать функцию:
\[
\min_x f(x),\qquad x=(x_1,x_2,\ldots,x_n).
\]
В координатном спуске по очереди выбирают направления:
\[
e_1=(1,0,\ldots,0),\quad e_2=(0,1,\ldots,0),\quad \ldots,\quad e_n=(0,0,\ldots,1).
\]
На каждой координате решается одномерная задача:
\[
\alpha_k=\arg\min_\alpha f(x_k+\alpha e_i).
\]
После этого точка обновляется:
\[
x_{k+1}=x_k+\alpha_ke_i.
\]
Преимущество метода - простота. Недостаток - если линии уровня функции сильно вытянуты и повернуты относительно координатных осей, метод может сходиться медленно.
\textbf{Метод Пауэлла} - это метод прямого поиска, который не требует вычисления производных. Он использует набор направлений и последовательно минимизирует функцию вдоль каждого из них.
Идея метода Пауэлла:
\noindent\textbf{1.}\ Задать начальную точку $x_0$ и набор направлений $d_1,\ldots,d_n$.\par
\noindent\textbf{2.}\ Последовательно выполнить одномерную минимизацию вдоль каждого направления:\par
\[
x^{(i)}=x^{(i-1)}+\alpha_id_i.
\]
\noindent\textbf{3.}\ Построить новое направление:\par
\[
d_{\text{new}}=x^{(n)}-x^{(0)}.
\]
\noindent\textbf{4.}\ Одно из старых направлений заменить новым.\par
Метод Пауэлла пытается строить направления, которые лучше соответствуют форме функции. Для квадратичных функций при точном одномерном поиске он может быть очень эффективным.
\textbf{Метод Хука-Дживса} - это метод поиска без производных. Он состоит из двух типов шагов: исследующего поиска и поиска по образцу.
\textbf{Исследующий поиск} - вокруг текущей точки пробуют сдвиги по координатам:
\[
x_i+h_i,\qquad x_i-h_i.
\]
Если значение функции уменьшается, шаг принимается. Если улучшения нет, шаг по этой координате уменьшается.
\textbf{Поиск по образцу} - если исследующий поиск нашел новую лучшую точку, метод пытается продолжить движение в том же направлении:
\[
x_{\text{pattern}}=x_{\text{new}}+(x_{\text{new}}-x_{\text{old}}).
\]
Так метод использует уже найденное удачное направление.

\subsection{Сравнение методов}
Координатный спуск прост, но может быть медленным.
Метод Пауэлла строит новые направления и обычно работает лучше на гладких функциях.
Метод Хука-Дживса удобен, когда производные неизвестны или функция задана как черный ящик.

\subsection{Пример}
Пусть
\[
f(x,y)=(x-2)^2+(y+1)^2.
\]
Начнем с точки $(0,0)$.
При координатной оптимизации сначала минимизируем по $x$ при фиксированном $y=0$:
\[
f(x,0)=(x-2)^2+1.
\]
Минимум достигается при $x=2$. Получаем точку $(2,0)$.
Затем минимизируем по $y$:
\[
f(2,y)=(y+1)^2.
\]
Минимум достигается при $y=-1$. Получаем точку $(2,-1)$, то есть минимум найден.

\section{13. Метод Нельдер-Мида}
\textbf{Метод Нельдера-Мида} - это метод безусловной оптимизации, который не требует вычисления производных. Он относится к методам прямого поиска и использует геометрическую фигуру, называемую симплексом.
В пространстве размерности $n$ симплекс состоит из $n+1$ точки. Например:
\noindent\textbullet\ в одномерном случае это два конца отрезка;\par
\noindent\textbullet\ в двумерном случае это треугольник;\par
\noindent\textbullet\ в трехмерном случае это тетраэдр.\par
На каждой итерации метод сравнивает значения функции в вершинах симплекса и постепенно перемещает симплекс в сторону меньших значений функции.
Пусть вершины симплекса:
\[
x_1,x_2,\ldots,x_{n+1}.
\]
Их упорядочивают так, чтобы:
\[
f(x_1)\le f(x_2)\le \ldots \le f(x_{n+1}).
\]
Здесь $x_1$ - лучшая вершина, а $x_{n+1}$ - худшая.
Далее вычисляют центр тяжести всех точек, кроме худшей:
\[
x_c=\frac{1}{n}\sum_{i=1}^{n}x_i.
\]
\textbf{Отражение} - худшую точку отражают относительно центра:
\[
x_r=x_c+\alpha(x_c-x_{n+1}),
\]
обычно $\alpha=1$.
Если отраженная точка хорошая, ее принимают вместо худшей.
\textbf{Растяжение} - если отраженная точка оказалась лучше всех, пробуют пойти еще дальше:
\[
x_e=x_c+\gamma(x_r-x_c),
\]
где $\gamma>1$
\textbf{Сжатие} - если отражение не дало хорошего результата, симплекс сжимают:
\[
x_s=x_c+\beta(x_{n+1}-x_c),
\]
где $0<\beta<1$.
\textbf{Редукция} - если даже сжатие не помогает, все точки приближают к лучшей:
\[
x_i=x_1+\delta(x_i-x_1),\qquad i=2,\ldots,n+1,
\]
где $0<\delta<1$.
\subsection{Достоинства и недостатки}
Достоинства:
\noindent\textbullet\ не нужны производные;\par
\noindent\textbullet\ метод прост в реализации;\par
\noindent\textbullet\ хорошо подходит для задач небольшой размерности.\par

Недостатки:
\noindent\textbullet\ нет такой строгой гарантии быстрой сходимости, как у ньютоновских методов;\par
\noindent\textbullet\ плохо работает в больших размерностях;\par
\noindent\textbullet\ может застревать на сложных функциях.\par

\subsection{Пример}
Пусть нужно минимизировать функцию:
\[
f(x,y)=x^2+y^2.
\]
Начальный симплекс в двумерном случае можно взять как треугольник с вершинами:
\[
(1,0),\quad (0,1),\quad (1,1).
\]
Худшая вершина - $(1,1)$, потому что $f(1,1)=2$.
Центр двух остальных точек:
\[
x_c=\left(\frac12,\frac12\right).
\]
Отражение худшей точки дает:
\[
x_r=x_c+(x_c-(1,1))=(0,0).
\]
В точке $(0,0)$ значение функции минимально: $f(0,0)=0$. Поэтому метод сразу нашел хорошую точку.
\section{14. Метод кросс-энтропии и алгоритм естественной эволюции}
\textbf{Метод кросс-энтропии} - это вероятностный метод оптимизации. Вместо того чтобы двигать одну точку, он задает распределение вероятностей на пространстве решений и постепенно изменяет это распределение так, чтобы оно чаще порождало хорошие решения.
Пусть нужно минимизировать функцию:
\[
\min_x f(x).
\]
Выбирается параметризованное распределение $p(x;\theta)$, например нормальное распределение:
\[
x\sim N(\mu,\sigma^2).
\]
Здесь параметры $\theta$ - это среднее $\mu$ и дисперсия $\sigma^2$.
Общий алгоритм метода кросс-энтропии:
\noindent\textbf{1.}\ Задать начальные параметры распределения $\theta_0$.\par
\noindent\textbf{2.}\ Сгенерировать набор случайных точек $x_1,\ldots,x_N$.\par
\noindent\textbf{3.}\ Вычислить значения функции $f(x_i)$.\par
\noindent\textbf{4.}\ Выбрать несколько лучших точек, которые называют элитной выборкой.\par
\noindent\textbf{5.}\ Обновить параметры распределения так, чтобы новые точки чаще попадали в область элитных решений.\par
Для нормального распределения обновление простое:
\[
\mu_{\text{new}}=\frac{1}{m}\sum_{i=1}^{m}x_i^{elite},
\]
\[
\sigma_{\text{new}}^2=\frac{1}{m}\sum_{i=1}^{m}(x_i^{elite}-\mu_{\text{new}})^2.
\]
То есть среднее сдвигается к лучшим точкам, а разброс постепенно уменьшается.
Название связано с тем, что параметры распределения выбираются через минимизацию кросс-энтропии между идеальным распределением, сосредоточенным на хороших решениях, и текущей моделью распределения.
\textbf{Алгоритм естественной эволюции} - это также популяционный вероятностный метод. Он оптимизирует не одну точку, а параметры распределения, из которого генерируются кандидаты.
Обычно используется распределение:
\[
x\sim p(x;\theta).
\]
Цель - подобрать $\theta$, чтобы математическое ожидание качества было наилучшим:
\[
J(\theta)=E_{x\sim p(x;\theta)}[F(x)].
\]
Если задача максимизации, то $F(x)$ - функция приспособленности. Если задача минимизации, можно взять $F(x)=-f(x)$.
В методах естественной эволюции параметры обновляются в направлении \textbf{естественного градиента}:
\[
\theta_{k+1}=\theta_k+\alpha \widetilde{\nabla}_\theta J(\theta_k).
\]
Естественный градиент учитывает геометрию пространства распределений и часто дает более устойчивое обновление, чем обычный градиент.
\subsection{Сравнение}
Метод кросс-энтропии обычно проще: он выбирает элитные точки и подгоняет распределение под них.
Алгоритмы естественной эволюции используют градиентный подход по параметрам распределения и могут более плавно учитывать качество всех сгенерированных решений.
\subsection{Пример}
Пусть нужно минимизировать:
\[
f(x)=(x-3)^2.
\]
Начнем с распределения $x\sim N(0,10)$.
Метод генерирует много точек. Лучшие точки будут находиться ближе к $x=3$, потому что там значение функции меньше.
После отбора элитных точек среднее $\mu$ сдвинется от $0$ к $3$, а дисперсия уменьшится. Через несколько итераций распределение станет сосредоточено около минимума.

\section{15. Генетические алгоритмы для оптимизации}
\textbf{Генетический алгоритм} - это популяционный метод оптимизации, вдохновленный идеями естественного отбора. Он работает не с одной точкой, а с множеством кандидатов, которое называется популяцией.

Каждый кандидат называется \textbf{особью} или \textbf{хромосомой}. Хромосома кодирует возможное решение задачи. Например, решение может быть записано:
\noindent\textbullet\ двоичной строкой;\par
\noindent\textbullet\ вектором чисел;\par
\noindent\textbullet\ перестановкой элементов.\par

Для каждой особи вычисляется \textbf{функция приспособленности}. Она показывает, насколько хорошим является решение. Если решается задача минимизации $f(x)$, то приспособленность можно задать как:
\[
Fitness(x)=\frac{1}{1+f(x)}.
\]
Чем меньше $f(x)$, тем больше значение $Fitness(x)$.

Основные операции генетического алгоритма:

\textbf{Отбор} - выбор лучших особей для размножения. Более приспособленные особи имеют большую вероятность попасть в следующее поколение. Популярные варианты: рулеточный отбор, турнирный отбор, отбор лучших.

\textbf{Скрещивание} - получение новых решений из двух родительских. Например, для двоичных строк:
\[
101|011,\qquad 110|100
\]
после одноточечного скрещивания могут получиться:
\[
101100,\qquad 110011.
\]

\textbf{Мутация} - случайное изменение небольшой части хромосомы. Например, в двоичной строке один бит может измениться с $0$ на $1$. Мутация нужна, чтобы алгоритм не терял разнообразие и мог выходить из локальных минимумов.

\textbf{Элитарность} - перенос нескольких лучших особей в следующее поколение без изменений. Это защищает уже найденные хорошие решения от случайной потери.

Общий алгоритм:
\noindent\textbf{1.}\ Сгенерировать начальную популяцию.\par
\noindent\textbf{2.}\ Вычислить приспособленность всех особей.\par
\noindent\textbf{3.}\ Выбрать родителей.\par
\noindent\textbf{4.}\ Выполнить скрещивание и мутацию.\par
\noindent\textbf{5.}\ Сформировать новое поколение.\par
\noindent\textbf{6.}\ Повторять процесс до выполнения критерия остановки.\par

Критерием остановки может быть достижение нужного качества, ограничение числа поколений или отсутствие улучшения в течение нескольких итераций.

\subsection{Достоинства и недостатки}
Достоинства:
\noindent\textbullet\ не нужны производные;\par
\noindent\textbullet\ можно применять к дискретным, непрерывным и смешанным задачам;\par
\noindent\textbullet\ метод хорошо подходит для сложных функций с большим числом локальных минимумов;\par
\noindent\textbullet\ популяция позволяет исследовать несколько областей пространства одновременно.\par

Недостатки:
\noindent\textbullet\ нет гарантии нахождения точного глобального оптимума;\par
\noindent\textbullet\ нужно подбирать параметры: размер популяции, вероятность мутации, вероятность скрещивания;\par
\noindent\textbullet\ метод может требовать много вычислений функции.\par

\subsection{Пример}
Пусть нужно найти максимум функции:
\[
F(x)=x^2
\]
на множестве целых чисел от $0$ до $7$.
Число можно закодировать тремя битами:
\[
0=000,\quad 1=001,\quad \ldots,\quad 7=111.
\]
Начальная популяция может быть:
\[
001,\quad 011,\quad 100,\quad 110.
\]
Это соответствует числам $1,3,4,6$.
Значения функции:
\[
1,\quad 9,\quad 16,\quad 36.
\]
Лучшей особью является $110$, то есть число $6$. После отбора, скрещивания и мутации в популяции может появиться $111$, то есть число $7$. Для него:
\[
F(7)=49.
\]
Это оптимальное решение на заданном множестве.

\section{16. Стайные алгоритмы: алгоритм стаи частиц, светлячковый алгоритм, алгоритм бактериального поиска пищи}
\textbf{Стайные алгоритмы} - это популяционные методы оптимизации, которые используют поведение группы простых агентов. Каждый агент сам по себе действует просто, но вся группа может находить хорошие решения за счет обмена информацией и случайного поиска.
Обычно такие методы применяют, когда функция сложная, негладкая, имеет много локальных минимумов или производные неизвестны.
\subsection{Алгоритм стаи частиц}
\textbf{Алгоритм стаи частиц}, или PSO, моделирует движение частиц в пространстве решений. Каждая частица имеет:
\noindent\textbullet\ текущую позицию $x_i^k$;\par
\noindent\textbullet\ скорость $v_i^k$;\par
\noindent\textbullet\ лучшую найденную ею позицию $p_i^k$;\par
\noindent\textbullet\ лучшую позицию всей стаи $g^k$.\par
Позиция частицы - это возможное решение задачи. Для задачи минимизации качество позиции оценивается значением функции $f(x)$.
Обновление скорости:
\[
v_i^{k+1}=\omega v_i^k+c_1r_1(p_i^k-x_i^k)+c_2r_2(g^k-x_i^k).
\]
Обновление позиции:
\[
x_i^{k+1}=x_i^k+v_i^{k+1}.
\]
Здесь $\omega$ - коэффициент инерции, $c_1$ - вес личного опыта частицы, $c_2$ - вес опыта всей стаи, $r_1,r_2$ - случайные числа из $[0,1]$.
Общий алгоритм PSO:
\noindent\textbf{1.}\ Случайно сгенерировать начальные позиции и скорости частиц.\par
\noindent\textbf{2.}\ Для каждой частицы вычислить $f(x_i)$.\par
\noindent\textbf{3.}\ Запомнить личные лучшие точки $p_i$ и глобальную лучшую точку $g$.\par
\noindent\textbf{4.}\ Обновить скорости и позиции частиц.\par
\noindent\textbf{5.}\ Пересчитать значения функции и обновить $p_i$ и $g$.\par
\noindent\textbf{6.}\ Повторять до остановки.\par
\subsection{Пример PSO}
Пусть нужно минимизировать:
\[
f(x)=(x-3)^2.
\]
Рассмотрим одну частицу: $x_0=0$, $v_0=0$, ее личная лучшая точка $p_0=0$, а лучшая точка стаи $g=4$. Возьмем $\omega=0.5$, $c_1=1$, $c_2=1$, $r_1=0.3$, $r_2=0.5$.
Тогда:
\[
v_1=0.5\cdot 0+1\cdot 0.3(0-0)+1\cdot 0.5(4-0)=2.
\]
Новая позиция:
\[
x_1=x_0+v_1=0+2=2.
\]
Частица сдвинулась ближе к минимуму $x=3$.
\subsection{Светлячковый алгоритм}
\textbf{Светлячковый алгоритм} основан на идее, что более яркие светлячки притягивают менее ярких. В оптимизации яркость связана со значением целевой функции. Для минимизации более ярким можно считать решение с меньшим $f(x)$.
Если светлячок $j$ лучше светлячка $i$, то $i$ движется к $j$:
\[
x_i^{k+1}=x_i^k+\beta_0e^{-\gamma r_{ij}^2}(x_j^k-x_i^k)+\alpha \varepsilon_i.
\]
Здесь $r_{ij}=\|x_i-x_j\|$ - расстояние между светлячками, $\beta_0$ - начальная привлекательность, $\gamma$ - коэффициент поглощения света, $\alpha \varepsilon_i$ - случайное возмущение.
Общий алгоритм:
\noindent\textbf{1.}\ Сгенерировать начальную популяцию светлячков.\par
\noindent\textbf{2.}\ Вычислить яркость каждого светлячка по значению функции.\par
\noindent\textbf{3.}\ Для каждой пары светлячков переместить худшего к лучшему.\par
\noindent\textbf{4.}\ Добавить случайное движение.\par
\noindent\textbf{5.}\ Обновить лучшие найденные решения.\par
\noindent\textbf{6.}\ Повторять до остановки.\par
\subsection{Алгоритм бактериального поиска пищи}
\textbf{Алгоритм бактериального поиска пищи} моделирует движение бактерий в поиске питательных веществ. Основные операции:
\noindent\textbullet\ хемотаксис - бактерия случайно меняет направление и движется;\par
\noindent\textbullet\ плавание - если движение улучшает значение функции, бактерия продолжает двигаться в том же направлении;\par
\noindent\textbullet\ размножение - лучшие бактерии копируются, худшие удаляются;\par
\noindent\textbullet\ ликвидация и рассеивание - часть бактерий случайно переносится в новые точки.\par
Шаг хемотаксиса можно записать так:
\[
x_i^{k+1}=x_i^k+C_i\frac{\Delta_i}{\|\Delta_i\|},
\]
где $C_i$ - длина шага, $\Delta_i$ - случайное направление.
Смысл метода: бактерии исследуют пространство, хорошие области усиливаются размножением, а случайное рассеивание помогает выходить из локальных минимумов.
\subsection{Сравнение}
PSO хорошо использует общий опыт стаи и часто быстро находит хорошие области.
Светлячковый алгоритм удобен для мультимодальных задач, потому что разные группы светлячков могут двигаться к разным локальным минимумам.
Бактериальный поиск пищи сочетает локальное улучшение и случайные скачки, но имеет много параметров.

\section{17. Условия Каруша-Куна-Такера}
\textbf{Условия Каруша-Куна-Такера}, или KKT, - это необходимые условия оптимальности для задач с ограничениями. Для выпуклых задач при выполнении дополнительных условий они становятся еще и достаточными.
Рассмотрим задачу:
\[
\min_x f(x)
\]
при ограничениях:
\[
g_i(x)\le 0,\qquad i=1,\ldots,m,
\]
\[
h_j(x)=0,\qquad j=1,\ldots,p.
\]
Строится функция Лагранжа:
\[
L(x,\lambda,\mu)=f(x)+\sum_{i=1}^m\lambda_i g_i(x)+\sum_{j=1}^p\mu_j h_j(x).
\]
Здесь $\lambda_i$ - множители для неравенств, а $\mu_j$ - множители для равенств.
Условия KKT:
\noindent\textbf{1.}\ \textbf{Допустимость прямой задачи}:\par
\[
g_i(x^*)\le 0,\qquad h_j(x^*)=0.
\]
\noindent\textbf{2.}\ \textbf{Двойственная допустимость}:\par
\[
\lambda_i\ge 0.
\]
\noindent\textbf{3.}\ \textbf{Стационарность}:\par
\[
\nabla_x L(x^*,\lambda^*,\mu^*)=0.
\]
То есть:
\[
\nabla f(x^*)+\sum_{i=1}^m\lambda_i^*\nabla g_i(x^*)+\sum_{j=1}^p\mu_j^*\nabla h_j(x^*)=0.
\]
\noindent\textbf{4.}\ \textbf{Дополняющая нежесткость}:\par
\[
\lambda_i^*g_i(x^*)=0,\qquad i=1,\ldots,m.
\]
Последнее условие означает: если ограничение неактивно, то его множитель равен нулю. Если множитель положительный, то ограничение выполняется как равенство.
\subsection{Пример}
Пусть нужно решить:
\[
\min_x x^2
\]
при ограничении:
\[
x\ge 1.
\]
Перепишем ограничение в стандартном виде:
\[
g(x)=1-x\le 0.
\]
Лагранжиан:
\[
L(x,\lambda)=x^2+\lambda(1-x).
\]
Условия KKT:
\[
2x-\lambda=0,
\]
\[
1-x\le 0,\qquad \lambda\ge 0,
\]
\[
\lambda(1-x)=0.
\]
Так как минимум без ограничения был бы в точке $x=0$, но она недопустима, оптимум находится на границе $x=1$.
Тогда из стационарности:
\[
2\cdot 1-\lambda=0.
\]
Получаем:
\[
\lambda=2.
\]
Все условия KKT выполнены, значит $x^*=1$ - оптимальное решение.

\section{18. Методы штрафных и барьерных функций}
\textbf{Методы штрафных и барьерных функций} позволяют свести задачу с ограничениями к последовательности задач без ограничений или с более простыми ограничениями.
Пусть есть задача:
\[
\min_x f(x)
\]
при ограничениях:
\[
g_i(x)\le 0,\qquad h_j(x)=0.
\]
\subsection{Метод штрафных функций}
В методе штрафов к целевой функции добавляют слагаемые, которые сильно увеличивают значение функции при нарушении ограничений.
Один из вариантов штрафной функции:
\[
F_r(x)=f(x)+r\sum_i\max(0,g_i(x))^2+r\sum_j h_j(x)^2.
\]
Здесь $r>0$ - параметр штрафа. Чем больше $r$, тем сильнее наказывается нарушение ограничений.
Общий алгоритм:
\noindent\textbf{1.}\ Выбрать начальное значение $r$.\par
\noindent\textbf{2.}\ Минимизировать функцию $F_r(x)$ без ограничений.\par
\noindent\textbf{3.}\ Увеличить $r$.\par
\noindent\textbf{4.}\ Снова минимизировать новую штрафную функцию.\par
\noindent\textbf{5.}\ Повторять, пока ограничения не будут выполняться с нужной точностью.\par
Преимущество метода - простая идея. Недостаток - при больших $r$ задача может становиться плохо обусловленной.
\subsection{Метод барьерных функций}
В методе барьеров, наоборот, точка должна оставаться внутри допустимой области. При приближении к границе функция резко возрастает.
Для ограничений $g_i(x)<0$ часто используют логарифмический барьер:
\[
F_\mu(x)=f(x)-\mu\sum_i\ln(-g_i(x)).
\]
Здесь $\mu>0$ - параметр барьера. Когда $x$ приближается к границе $g_i(x)=0$, выражение $\ln(-g_i(x))$ стремится к $-\infty$, поэтому $F_\mu(x)$ стремится к $+\infty$.
Общий алгоритм:
\noindent\textbf{1.}\ Выбрать допустимую внутреннюю точку $x_0$.\par
\noindent\textbf{2.}\ Выбрать начальное $\mu>0$.\par
\noindent\textbf{3.}\ Минимизировать барьерную функцию $F_\mu(x)$.\par
\noindent\textbf{4.}\ Уменьшить $\mu$.\par
\noindent\textbf{5.}\ Повторять, постепенно приближаясь к решению исходной задачи.\par
\subsection{Отличие штрафа от барьера}
Штрафной метод может рассматривать недопустимые точки, но наказывает их.
Барьерный метод не выходит за границу допустимой области и требует начальную внутреннюю точку.
В штрафном методе параметр $r$ обычно увеличивают, а в барьерном методе параметр $\mu$ уменьшают.
\subsection{Пример}
Пусть нужно минимизировать:
\[
f(x)=(x-2)^2
\]
при ограничении:
\[
x\le 1.
\]
Истинное решение находится на границе: $x^*=1$.
Для штрафного метода можно взять:
\[
F_r(x)=(x-2)^2+r\max(0,x-1)^2.
\]
Если $x>1$, включается штраф за нарушение ограничения. При росте $r$ минимум $F_r(x)$ будет приближаться к $x=1$.
Для барьерного метода берем внутреннюю область $x<1$:
\[
F_\mu(x)=(x-2)^2-\mu\ln(1-x).
\]
При $x\to 1-0$ барьер стремится к бесконечности, поэтому метод подходит к границе только изнутри.

\section{19. Метод Зойтендейка}
\textbf{Метод Зойтендейка} - это метод допустимых направлений для задач условной оптимизации. Он строит направление, которое одновременно уменьшает целевую функцию и не нарушает активные ограничения в первом приближении.
Рассмотрим задачу:
\[
\min_x f(x)
\]
при ограничениях:
\[
g_i(x)\le 0,\qquad i=1,\ldots,m.
\]
На текущей допустимой точке $x_k$ выделяют активные ограничения:
\[
I(x_k)=\{i:\ g_i(x_k)=0\}.
\]
Направление $d_k$ должно быть направлением спуска:
\[
\nabla f(x_k)^T d_k<0,
\]
и не должно сразу выводить точку из допустимой области по активным ограничениям:
\[
\nabla g_i(x_k)^Td_k\le 0,\qquad i\in I(x_k).
\]
Если такое направление найдено, новая точка ищется в виде:
\[
x_{k+1}=x_k+\alpha_k d_k,
\]
где $\alpha_k>0$ выбирается так, чтобы точка оставалась допустимой и значение функции уменьшалось.
\subsection{Как ищут направление}
В методе Зойтендейка направление часто находят как решение вспомогательной задачи:
\[
\min_{d,\theta}\ \theta
\]
при условиях:
\[
\nabla f(x_k)^Td\le \theta,
\]
\[
\nabla g_i(x_k)^Td\le \theta,\qquad i\in I(x_k),
\]
\[
\|d\|\le 1.
\]
Если получилось $\theta<0$, то найдено допустимое направление убывания. Если $\theta\ge 0$, то допустимого направления убывания нет, и точка считается кандидатом на оптимум.
\subsection{Общий алгоритм}
\noindent\textbf{1.}\ Выбрать начальную допустимую точку $x_0$.\par
\noindent\textbf{2.}\ Найти активные ограничения.\par
\noindent\textbf{3.}\ Решить вспомогательную задачу для направления $d_k$.\par
\noindent\textbf{4.}\ Если направления убывания нет, остановиться.\par
\noindent\textbf{5.}\ Выбрать шаг $\alpha_k$ с сохранением допустимости.\par
\noindent\textbf{6.}\ Обновить точку $x_{k+1}=x_k+\alpha_kd_k$.\par
\noindent\textbf{7.}\ Повторять до выполнения критерия остановки.\par
\subsection{Пример}
Пусть нужно минимизировать:
\[
f(x,y)=(x-2)^2+y^2
\]
при ограничении:
\[
x+y\le 1.
\]
В точке $(0,0)$ ограничение неактивно, потому что $0+0<1$. Поэтому можно идти по антиградиенту. Градиент:
\[
\nabla f(0,0)=(-4,0).
\]
Антиградиент:
\[
d=(4,0).
\]
Это направление уменьшает функцию, но слишком большой шаг может вывести за границу $x+y\le 1$.
Если точка уже на границе, например $x_k=(1,0)$, то активное ограничение имеет градиент:
\[
\nabla g(x_k)=(1,1).
\]
Допустимое направление должно удовлетворять:
\[
d_x+d_y\le 0.
\]
Например, направление $d=(0.5,-1)$ допустимо, потому что $0.5-1=-0.5\le 0$. При этом:
\[
\nabla f(1,0)=(-2,0),
\]
\[
\nabla f(1,0)^Td=-1<0.
\]
Значит, это допустимое направление убывания.

\section{20. Метод проекции градиента}
\textbf{Метод проекции градиента} применяется для задач с ограничением на множество допустимых точек:
\[
\min_{x\in \Omega} f(x).
\]
Обычный градиентный шаг может вывести точку за пределы множества $\Omega$. Поэтому после шага точку проецируют обратно на допустимое множество.
Итерация метода:
\[
y_k=x_k-\alpha_k\nabla f(x_k),
\]
\[
x_{k+1}=P_\Omega(y_k).
\]
Здесь $P_\Omega(y)$ - проекция точки $y$ на множество $\Omega$:
\[
P_\Omega(y)=\arg\min_{z\in \Omega}\|z-y\|.
\]
То есть выбирается ближайшая допустимая точка.
\subsection{Общий алгоритм}
\noindent\textbf{1.}\ Выбрать допустимую начальную точку $x_0\in \Omega$.\par
\noindent\textbf{2.}\ Сделать градиентный шаг без учета ограничений.\par
\noindent\textbf{3.}\ Спроецировать полученную точку на множество $\Omega$.\par
\noindent\textbf{4.}\ Повторять, пока точка или значение функции почти не меняются.\par
Для простого отрезка $\Omega=[a,b]$ проекция имеет вид:
\[
P_{[a,b]}(y)=\min(\max(y,a),b).
\]
То есть если $y<a$, берется $a$; если $y>b$, берется $b$; если $y\in [a,b]$, точка не меняется.
\subsection{Пример}
Пусть нужно минимизировать:
\[
f(x)=(x-3)^2
\]
на отрезке:
\[
0\le x\le 2.
\]
Без ограничения минимум был бы в точке $x=3$, но она недопустима. Значит, минимум на отрезке будет в точке $x=2$.
Возьмем $x_0=0$ и $\alpha=1$. Градиент:
\[
f'(x)=2(x-3),
\]
\[
f'(0)=-6.
\]
Обычный градиентный шаг:
\[
y_0=x_0-\alpha f'(x_0)=0-1\cdot(-6)=6.
\]
Но $6$ не принадлежит отрезку $[0,2]$. Проецируем:
\[
x_1=P_{[0,2]}(6)=2.
\]
После одного шага метод пришел в оптимальную допустимую точку.

\section{21. Последовательное квадратичное программирование}
\textbf{Последовательное квадратичное программирование}, или SQP, - это метод решения нелинейных задач оптимизации с ограничениями. Его идея: на каждой итерации заменить исходную нелинейную задачу квадратичной подзадачей.
Рассмотрим задачу:
\[
\min_x f(x)
\]
при ограничениях:
\[
g_i(x)\le 0,\qquad h_j(x)=0.
\]
В точке $x_k$ строится квадратичная модель целевой функции и линейные приближения ограничений.
Квадратичная подзадача имеет вид:
\[
\min_p\ \nabla f(x_k)^Tp+\frac12p^TB_kp
\]
при условиях:
\[
g_i(x_k)+\nabla g_i(x_k)^Tp\le 0,
\]
\[
h_j(x_k)+\nabla h_j(x_k)^Tp=0.
\]
Здесь $p$ - шаг, а $B_k$ - матрица Гессе лагранжиана или ее квазиньютоновское приближение.
После решения квадратичной подзадачи получают направление $p_k$ и обновляют точку:
\[
x_{k+1}=x_k+\alpha_kp_k.
\]
Шаг $\alpha_k$ часто выбирают с помощью одномерного поиска по штрафной или merit-функции, чтобы одновременно уменьшать целевую функцию и нарушение ограничений.
\subsection{Общий алгоритм SQP}
\noindent\textbf{1.}\ Выбрать начальную точку $x_0$ и матрицу $B_0$.\par
\noindent\textbf{2.}\ Построить квадратичную подзадачу в точке $x_k$.\par
\noindent\textbf{3.}\ Решить эту подзадачу и получить направление $p_k$.\par
\noindent\textbf{4.}\ Выбрать шаг $\alpha_k$.\par
\noindent\textbf{5.}\ Обновить точку $x_{k+1}=x_k+\alpha_kp_k$.\par
\noindent\textbf{6.}\ Обновить множители Лагранжа и матрицу $B_k$.\par
\noindent\textbf{7.}\ Повторять до выполнения условий KKT с нужной точностью.\par
SQP можно понимать как метод Ньютона, примененный к условиям KKT. Поэтому около решения он может сходиться очень быстро.
\subsection{Пример}
Пусть нужно минимизировать:
\[
f(x,y)=(x-1)^2+(y-2)^2
\]
при ограничении:
\[
x+y=1.
\]
Возьмем начальную точку $x_0=(0,0)^T$. Обозначим шаг:
\[
p=(p_x,p_y)^T.
\]
Градиент целевой функции в точке $(0,0)$:
\[
\nabla f(0,0)=(-2,-4)^T.
\]
Для этой квадратичной функции можно взять $B=2I$. Линеаризованное ограничение:
\[
0+0+p_x+p_y=1,
\]
то есть:
\[
p_x+p_y=1.
\]
Квадратичная подзадача:
\[
\min_p\ -2p_x-4p_y+p_x^2+p_y^2
\]
при условии:
\[
p_x+p_y=1.
\]
Решение этой подзадачи:
\[
p=(0,1)^T.
\]
Тогда:
\[
(x_1,y_1)=(0,0)+(0,1)=(0,1).
\]
Эта точка удовлетворяет ограничению и является точкой минимума исходной задачи.

\section{22. Двойственная задача. Численные методы, основанные на двойственной задаче: Primal-Dual algorithm, Dual ascent}
\textbf{Двойственная задача} строится по исходной, или прямой, задаче оптимизации и часто позволяет получить нижние оценки оптимального значения, а иногда решить задачу удобнее.
Рассмотрим прямую задачу:
\[
\min_x f(x)
\]
при ограничениях:
\[
g_i(x)\le 0,\qquad h_j(x)=0.
\]
Лагранжиан:
\[
L(x,\lambda,\mu)=f(x)+\sum_i\lambda_i g_i(x)+\sum_j\mu_j h_j(x),
\]
где $\lambda_i\ge 0$.
Двойственная функция:
\[
q(\lambda,\mu)=\inf_x L(x,\lambda,\mu).
\]
Двойственная задача:
\[
\max_{\lambda,\mu}\ q(\lambda,\mu)
\]
при условии:
\[
\lambda_i\ge 0.
\]
Для задачи минимизации двойственная функция дает нижнюю оценку оптимального значения прямой задачи:
\[
q(\lambda,\mu)\le f(x^*).
\]
Это называется слабой двойственностью. Для выпуклых задач при выполнении условий регулярности может выполняться сильная двойственность: оптимальные значения прямой и двойственной задач совпадают.
\subsection{Dual ascent}
\textbf{Dual ascent} обновляет двойственные переменные в направлении роста двойственной функции. На каждой итерации:
\noindent\textbf{1.}\ При фиксированных $\lambda^k,\mu^k$ найти:\par
\[
x^k=\arg\min_x L(x,\lambda^k,\mu^k).
\]
\noindent\textbf{2.}\ Обновить множители:\par
\[
\lambda_i^{k+1}=\max(0,\lambda_i^k+\alpha_k g_i(x^k)),
\]
\[
\mu_j^{k+1}=\mu_j^k+\alpha_k h_j(x^k).
\]
Если ограничение $g_i(x)\le 0$ нарушено, то $g_i(x)>0$, и множитель $\lambda_i$ увеличивается.
\subsection{Primal-Dual algorithm}
\textbf{Primal-Dual algorithm} одновременно обновляет прямые переменные $x$ и двойственные переменные $\lambda,\mu$.
Типичная схема:
\[
x^{k+1}=x^k-\alpha_k\nabla_x L(x^k,\lambda^k,\mu^k),
\]
\[
\lambda_i^{k+1}=\max(0,\lambda_i^k+\alpha_k g_i(x^k)),
\]
\[
\mu_j^{k+1}=\mu_j^k+\alpha_k h_j(x^k).
\]
То есть по $x$ делают шаг спуска, а по двойственным переменным - шаг подъема.
\subsection{Пример двойственной задачи}
Пусть:
\[
\min_x x^2
\]
при ограничении:
\[
x\ge 1.
\]
Запишем ограничение как:
\[
g(x)=1-x\le 0.
\]
Лагранжиан:
\[
L(x,\lambda)=x^2+\lambda(1-x).
\]
Чтобы найти двойственную функцию, минимизируем по $x$:
\[
\frac{dL}{dx}=2x-\lambda=0.
\]
Отсюда:
\[
x=\frac{\lambda}{2}.
\]
Подставляем:
\[
q(\lambda)=\left(\frac{\lambda}{2}\right)^2+\lambda\left(1-\frac{\lambda}{2}\right)
=\lambda-\frac{\lambda^2}{4}.
\]
Двойственная задача:
\[
\max_{\lambda\ge 0}\left(\lambda-\frac{\lambda^2}{4}\right).
\]
Производная:
\[
1-\frac{\lambda}{2}=0.
\]
Получаем:
\[
\lambda^*=2.
\]
Тогда:
\[
x^*=\frac{\lambda^*}{2}=1.
\]
Это совпадает с решением прямой задачи.

\section{23. Многокритериальная оптимизация, оптимальность по Парето. Методы построения границы Парето: метод ограничений, метод взвешивания, обобщенный минимаксный метод}
\textbf{Многокритериальная оптимизация} - это оптимизация сразу нескольких целевых функций. Обычно улучшение одного критерия может ухудшать другой.
Задача записывается так:
\[
\min_x F(x)=\left(f_1(x),f_2(x),\ldots,f_m(x)\right).
\]
\textbf{Доминирование по Парето} - точка $x$ доминирует точку $y$, если:
\[
f_i(x)\le f_i(y)\quad \text{для всех } i,
\]
и хотя бы для одного критерия неравенство строгое:
\[
f_j(x)<f_j(y).
\]
\textbf{Оптимальность по Парето} - точка $x^*$ называется Парето-оптимальной, если не существует другой допустимой точки, которая улучшает хотя бы один критерий и не ухудшает остальные.
Множество Парето-оптимальных точек в пространстве переменных называют множеством Парето, а их образы в пространстве критериев - границей Парето.
\subsection{Метод взвешивания}
Несколько критериев заменяют одной функцией:
\[
\min_x \sum_{i=1}^m w_i f_i(x),
\]
где:
\[
w_i\ge 0,\qquad \sum_{i=1}^m w_i=1.
\]
Меняя веса $w_i$, получают разные точки границы Парето.
Плюс метода - простота. Минус - для невыпуклой границы Парето метод может пропускать некоторые участки.
\subsection{Метод ограничений}
Один критерий оставляют целевой функцией, а остальные переводят в ограничения:
\[
\min_x f_1(x)
\]
при условиях:
\[
f_i(x)\le \varepsilon_i,\qquad i=2,\ldots,m.
\]
Меняя уровни $\varepsilon_i$, получают разные компромиссные решения.
Этот метод может находить точки и на невыпуклой границе Парето, но требует подбора ограничений.
\subsection{Обобщенный минимаксный метод}
В этом методе минимизируют максимальное нормированное отклонение от идеальной точки. Пусть $z^*$ - идеальная точка, где каждая компонента равна лучшему возможному значению соответствующего критерия. Тогда решают задачу:
\[
\min_x \max_i w_i\left(f_i(x)-z_i^*\right).
\]
Смысл: найти решение, у которого наихудшее отклонение от идеала минимально.
\subsection{Пример}
Пусть:
\[
f_1(x)=x^2,\qquad f_2(x)=(x-2)^2,
\]
где:
\[
x\in[0,2].
\]
Первый критерий хочет $x=0$, второй хочет $x=2$. Поэтому все точки между $0$ и $2$ являются компромиссными.
Метод взвешивания дает:
\[
\min_x\ w x^2+(1-w)(x-2)^2.
\]
Производная:
\[
2wx+2(1-w)(x-2)=0.
\]
Отсюда:
\[
x=2(1-w).
\]
При $w=1$ получаем $x=0$, при $w=0$ получаем $x=2$, а при промежуточных весах получаем промежуточные точки границы Парето.

\section{24. Популяционные методы нахождения границы Парето: nondomination ranking, SPEA, vector evaluated genetic algorithm. Фильтр Парето}
Популяционные методы удобны для многокритериальной оптимизации, потому что они работают сразу с множеством решений. Это позволяет строить не одну точку, а приближение всей границы Парето.
\subsection{Nondomination ranking}
\textbf{Nondomination ranking} - это ранжирование решений по уровням недоминирования.
Первый фронт $F_1$ состоит из решений, которые не доминируются никакими другими решениями популяции. Это лучшие решения по Парето.
После удаления $F_1$ среди оставшихся решений находят следующий недоминируемый фронт $F_2$. Затем аналогично находят $F_3$ и так далее.
Алгоритм:
\noindent\textbf{1.}\ Для каждого решения определить, какие решения оно доминирует и сколько решений доминируют его.\par
\noindent\textbf{2.}\ Все решения, которые никто не доминирует, поместить во фронт $F_1$.\par
\noindent\textbf{3.}\ Удалить влияние решений из $F_1$ и найти следующий фронт.\par
\noindent\textbf{4.}\ Повторять, пока все решения не получат ранг.\par
Чем меньше номер фронта, тем лучше решение. В генетических алгоритмах решения из лучших фронтов чаще выбираются для размножения.
\subsection{SPEA}
\textbf{SPEA}, или Strength Pareto Evolutionary Algorithm, - это эволюционный алгоритм для многокритериальной оптимизации. Он использует внешнее хранилище, или архив, куда помещаются недоминируемые решения.
Идея SPEA:
\noindent\textbf{1.}\ Сгенерировать популяцию.\par
\noindent\textbf{2.}\ Найти недоминируемые решения и добавить их в архив.\par
\noindent\textbf{3.}\ Каждому решению назначить силу, зависящую от того, сколько решений оно доминирует.\par
\noindent\textbf{4.}\ Рассчитать приспособленность решений с учетом доминирования и плотности.\par
\noindent\textbf{5.}\ Выбрать родителей из популяции и архива.\par
\noindent\textbf{6.}\ Выполнить скрещивание и мутацию.\par
\noindent\textbf{7.}\ Повторять, обновляя архив Парето-решений.\par
Архив нужен, чтобы хорошие недоминируемые решения не терялись при случайных изменениях популяции.
\subsection{Vector evaluated genetic algorithm}
\textbf{Vector Evaluated Genetic Algorithm}, или VEGA, - один из ранних многокритериальных генетических алгоритмов.
Главная идея: популяцию делят на группы по числу критериев. Каждая группа отбирается по одному критерию.
Например, если есть два критерия $f_1$ и $f_2$, то:
\noindent\textbullet\ первая часть популяции отбирается по качеству $f_1$;\par
\noindent\textbullet\ вторая часть популяции отбирается по качеству $f_2$.\par
Потом выбранные решения объединяются, и к ним применяются скрещивание и мутация.
Метод простой, но может давать перекос к крайним решениям, которые хороши по одному критерию и плохи по другому.
\subsection{Фильтр Парето}
\textbf{Фильтр Парето} - это процедура поддержания множества недоминируемых решений.
Когда появляется новое решение $x$, его сравнивают с решениями в текущем архиве:
\noindent\textbf{1.}\ Если новое решение доминируется хотя бы одним архивным решением, его отбрасывают.\par
\noindent\textbf{2.}\ Если новое решение доминирует какие-то архивные решения, эти решения удаляют.\par
\noindent\textbf{3.}\ Если новое решение не доминируется, его добавляют в архив.\par
В итоге в архиве остаются только недоминируемые решения.
\subsection{Пример фильтра Парето}
Пусть минимизируются два критерия, и есть точки в пространстве критериев:
\[
A=(1,5),\quad B=(2,3),\quad C=(4,2),\quad D=(3,4),\quad E=(5,1).
\]
Точка $B=(2,3)$ доминирует точку $D=(3,4)$, потому что:
\[
2<3,\qquad 3<4.
\]
Значит, $D$ удаляется из Парето-фильтра.
Точки $A,B,C,E$ не доминируют друг друга: каждая из них лучше по одному критерию и хуже по другому. Поэтому они остаются в архиве как приближение границы Парето.
\subsection{Сравнение}
Nondomination ranking удобно использовать для отбора в популяции.
SPEA дополнительно хранит внешний архив лучших Парето-решений.
VEGA прост, но может плохо сохранять баланс между критериями.
Фильтр Парето - базовая процедура, которая оставляет только недоминируемые решения.

\end{document}
